% concatenated from submersions.tex
%%%%%%%%%%%%%%%%%%%%%%%%%%%%%%%%%%%%%%%%%%%%%%%%%%%%%%%%%%%%%
%
% 
%      AMS-LaTeX file for 
%
%      
%          
%%%%%%%%%%%%%%%%%%%%%%%%%%%%%%%%%%%%%%%%%%%%%%%%%%%%%%%%%%%%%

\documentclass[11pt]{amsart}
\usepackage{amsmath,amscd,amssymb,amsfonts,amsthm,hyperref}
\usepackage{tcolorbox}
\usepackage{tikz-cd}
\hypersetup{colorlinks,   linkcolor={red!50!black},    citecolor={blue!70!black},   urlcolor={blue!80!black}}
\usepackage[cmtip,matrix,arrow]{xy}
\usepackage{mathabx}
\usepackage{graphicx}
\usepackage{tikz-cd}
\usepackage{MnSymbol}   


\textwidth15.8 cm
\oddsidemargin.4cm
\evensidemargin.4cm
% 
%TEMP

%\CompileMatrices

%\newdir{((}{{\lhook\kern-.1em}}
%\newdir{))}{{\rhook\kern-.1em}}
%\newdir{ >}{{}*!/-5pt/@{>}}
	
\renewcommand\theenumi{\alph{enumi}}
\renewcommand\theenumii{\roman{enumii}}

\newtheorem{theorem}{Theorem}[section]
\newtheorem{corollary}[theorem]{Corollary}
\newtheorem{proposition}[theorem]{Proposition}
\newtheorem{theo}{Theorem}[section]
\newtheorem{lemma}[theorem]{Lemma}
\theoremstyle{definition}    
\newtheorem{definition}[theorem]{Definition}
\theoremstyle{remark}
\newtheorem{assumption}[theorem]{Assumption}
\newtheorem{remark}[theorem]{Remark}
\newtheorem{remarks}[theorem]{Remarks}
\newtheorem{example}[theorem]{Example}
\newtheorem{examples}[theorem]{Examples}
%% Calligraphic and bold face letters 
%
\newcommand{\Ra}{\Rightarrow}
\newcommand\A{\mathcal{A}}
\newcommand{\EE}{\mathsf{E}}
\newcommand{\Cour}[1]      {[\![#1]\!]}
\newcommand\M{\mathcal{M}}
\newcommand\G{\mathcal{G}}
\newcommand{\K}{\mathbb{K}}
\newcommand{\W}{\mathcal{W}}
\newcommand{\V}{\mathcal{V}}
\renewcommand{\L}{\mathcal{L}}
\renewcommand{\O}{\mathcal{O}}
\newcommand{\T}{\mathbb{T}}
\newcommand{\J}{\mathcal{J}}
\newcommand{\Co}{\mathcal{C}}
\newcommand{\ca}{\mathcal}
\newcommand{\cg}{g}
\newcommand{\U}{\on{U}}
\newcommand{\E}{\ca{E}}
\newcommand{\X}{\mf{X}}
\newcommand{\F}{\mathcal{F}}
\newcommand{\N}{\mathbb{N}}
\newcommand{\R}{\mathbb{R}}
\newcommand{\C}{\mathbb{C}}
\newcommand{\SU}{\on{SU}}
\newcommand{\Y}{\mathcal{Y}}
\newcommand{\Z}{\mathbb{Z}}
\newcommand{\Q}{\ca{Q}}
\newcommand\pt{\on{pt}}
\renewcommand{\P}{\ca{P}}
% Lie algebras
\newcommand{\uu}{u}
\newcommand{\id}{\on{id}}

\newcommand\lie[1]{\mathfrak{#1}}
\renewcommand{\k}{\lie{k}}
\newcommand{\h}{\lie{h}}
\newcommand{\g}{\lie{g}}
\newcommand{\m}{\lie{m}}
\newcommand{\z}{\lie{z}}

\renewcommand{\t}{\lie{t}}
\newcommand{\Alc}{\lie{A}}

% Operatornames
\newcommand{\on}{\operatorname}
\newcommand{\Av}{ \on{Av} }
\newcommand{\length}{\on{length}}
\newcommand{\Eul}{ \on{Eul} } 
\newcommand{\Aut}{ \on{Aut} } 
\newcommand{\Ad}{ \on{Ad} }
\newcommand{\ad}{\on{ad}}
\newcommand{\Pf}{\on{Pf}}  % Pfaffian 
\newcommand{\Hol}{ \on{Hol} } 
\newcommand{\Rep}{\on{Rep}}
\newcommand{\Car}{\on{Car}}
\newcommand{\Hom}{ \on{Hom}}
\newcommand{\Ind}{ \on{Ind}}
\newcommand{\Irr}{ \on{Irr}}
\newcommand{\Res}{\on{Res}}
\renewcommand{\ker}{ \on{ker}}
\newcommand{\im}{ \on{im}}
\newcommand{\ind}{ \on{ind}}
\newcommand{\Spin}{ \on{Spin}}
\newcommand{\Pin}{ \on{Pin}}
\newcommand{\SO}{ \on{SO}}
\newcommand{\Mult}{  \on{Mult}}
\newcommand{\Vol}{  \on{Vol}}
\newcommand{\vol}{  \on{vol}}
\newcommand{\diag}{  \on{diag}}
\newcommand{\Pair}{\on{Pair}}
\renewcommand{\dh}{\m}
% Other macros
\renewcommand{\index}{ \on{index}}
\newcommand{\codim}{\on{codim}}
\newcommand{\cone}{ \on{cone} }
\newcommand{\tpi}{{2\pi i}}
%\newcommand{\tpi}{{2\pi\sqrt{-1}}}
\newcommand{\CS}{\on{CS}}
\newcommand{\da}{\dasharrow}

\newcommand\qu{/\kern-.7ex/} % Categorical quotients
\newcommand{\fus}{\circledast} % Fusion product
\newcommand{\Waff}{W_{\text{aff}}} % affine Weyl group 
\newcommand{\ipwc}{\on{int}(\t^\ast _+)} % interior of the positive Weyl chamber
\newcommand{\TG}{\mathbb{T}G}
\newcommand{\lra}{\longrightarrow}
\newcommand{\hra}{\hookrightarrow}
\newcommand{\ra}{\rightarrow}
\newcommand{\bib}{\bibitem}
\renewcommand{\d}{{\mbox{d}}}
\newcommand{\ol}{\overline}
\newcommand\Gpar{\G_\partial}
\newcommand\Phinv{\Phi^{-1}}
\newcommand\phinv{\phi^{-1}}
\newcommand\lam{\lambda}
\newcommand\Lam{\Lambda}
\newcommand\Sig{\Sigma}
\newcommand\sig{\sigma}
\newcommand\eps{\epsilon}
\newcommand\Om{\Omega}
\newcommand\om{\omega}
\newcommand{\Del}{\Delta}
\newcommand{\peps}{{(\eps)}}
\newcommand{\del}{\delta}
\newcommand{\f}{\frac}
\newcommand{\olt}{\overline{\theta}}
\newcommand{\p}{\partial}
\renewcommand{\l}{\langle}
\renewcommand{\r}{\rangle}
\newcommand\hh{{\f{1}{2}}}
\newcommand\tf{\tfrac}
\newcommand{\ti}{\tilde}
\newcommand{\eeq}{\end{eqnarray*}}
\newcommand{\beq}{\begin{eqnarray*}}
\newcommand{\hull}{\on{hull}}
\renewcommand{\H}{\ca{H}}
\newcommand{\La}{\Lambda}
\newcommand{\D}{\ca{D}}
\newcommand{\Cl}{\on{Cl}}
\newcommand{\pr}{\on{pr}}
\newcommand{\wh}{\widehat}
\newcommand{\wt}{\widetilde}
\newcommand{\mf}{\mathfrak}
\newcommand{\rra}{\rightrightarrows}
\newcommand{\sz}{\mathsf{s}}
\newcommand{\tz}{\mathsf{t}}
\newcommand{\iz}{\mathsf{i}}
\newcommand{\jz}{\mathsf{j}}
\newcommand{\pz}{\mathsf{p}}
\newcommand{\hz}{\mathsf{h}}


\newcommand{\core}{{\on{core}}}
\newcommand{\lin}{{\on{lin}}}
\renewcommand{\vert}{{\on{vert}}}
\newcommand{\tang}{{\on{tan}}}
\newcommand{\side}{{\on{side}}}
\newcommand{\LA}{\ca{LA}}
\newcommand{\VB}{\ca{VB}}
\newcommand{\CA}{\ca{CA}}
\newcommand{\LG}{\ca{LG}}
\newcommand{\gr}{\on{gr}}
\newcommand{\dd}{\mathfrak{d}}
\newcommand{\As}{A^{\l 2\r}}
\newcommand{\CC}{\mathsf{C}}
\newcommand{\BB}{\mathsf{B}}
\newcommand{\WW}{\mathsf{W}}
\renewcommand{\subset}{\subseteq}
\renewcommand{\supset}{\supseteq}

\setcounter{tocdepth}{1}



% *****************************************************************
% *****************************************************************
% *****************************************************************
% *****************************************************************
% *****************************************************************
% 
%       BEGINNING OF THE DOCUMENT
%
% *****************************************************************
% *****************************************************************
% *****************************************************************
\begin{document}
\sloppy
\title{A note on the cohomology of pullback Lie algebroids}
\author{Eckhard Meinrenken}

\begin{abstract}		
We give a direct differential-geometric proof of a result of Crainic on the cohomology of pullbacks of Lie algebroids under surjective submersions: 
If the fibers are connected and have vanishing cohomology through
degree $n$, pullback with coefficients in any representation is an
isomorphism through degree $n$ and is injective in degree $n+1$.
The proof reduces  to a  vanishing result for leafwise cohomology and uses an exhaustion-and-patching argument due to Buchdahl. No local triviality or finite-type assumption on the fibers is required. We also prove the analogous result for holomorphic Lie algebroids.
\end{abstract}

%\dedicatory{Dedicated to Rui Loja Fernandes on the occasion of his	sixtieth birthday.}
\maketitle

\section{Introduction}


Let $\pi\colon M\to N$ be a surjective submersion, and let
$A\Ra  N$ be a Lie algebroid. A fundamental result used in the proof of the van Est theorem for Lie groupoids \cite{cra:dif} 
asserts that, if the
fibers of $\pi$ are connected and have vanishing cohomology
in degrees $1,\ldots,n$, then pullback 
\[
\pi^*\colon H^k(A)\longrightarrow H^k(\pi^!A)
\]
is an isomorphism for $k\leq n$ and is injective in degree $n+1$.
More generally, the result holds with coefficients in a
representation.

This result was proved by Crainic~\cite[Theorem~2]{cra:dif} using a
sheaf-theoretic acyclicity criterion of Bernstein and
Lunts~\cite{ber:equ}. It is a key ingredient in his proof of
the van Est theorem, as well as in many subsequent variants of van Est
theory. The argument, however, is somewhat difficult to extract from
the cited sheaf-theoretic machinery, whose hypotheses and terminology
are not immediately transparent in the context of smooth
submersions.


The purpose of this note is to present a direct differential-geometric
proof. The main ingredient is an exhaustion-and-patching argument due
to Buchdahl \cite{buc:rel}, which we adapt to the present setting.
Our proof uses only standard tools from differential topology and de
Rham theory, at the level of Bott and Tu \cite{bo:di}. Besides making
the argument self-contained, this approach displays  the
main difficulties: a general surjective submersion need not be a locally
trivial fibration, and its fibers need not be of finite type. We also prove the corresponding result for holomorphic 
Lie algebroids and holomorphic submersions. 

\bigskip
\paragraph{\bf Acknowledgements.}
I thank Marius Crainic for discussions on the sidelines of the June
2025 Lisbon conference in honour of Rui Fernandes. I am also grateful to
Maarten Mol for patient explanations of the Bernstein--Lunts argument
and for pointing out Buchdahl's closely related result, which became
the starting point for this note.











\section{Pullback Lie algebroids}
%
All manifolds considered in this note are paracompact and Hausdorff. 
Let $A\Ra N$ be a Lie algebroid with anchor $\varrho\colon A\to TN$. The Lie algebroid cohomology of $A$ with coefficients in an $A$-representation 
$E\to N$ is denoted $H^\bullet(A,E)$. Given a smooth map $f\colon M\to N$ transverse to the anchor, one defines 
the pullback Lie algebroid $f^!A\Ra M$ (using standard notation as in \cite[Section 13.4]{cra:lec}) as the fiber product 
\[ f^!A=TM\times_{TN} A,\]
with Lie algebroid structure as a subalgebroid of $TM\times A$. Projection to the second factor gives a Lie algebroid morphism $f^!A\to A$. 
Given a representation of $A$ on a vector bundle $E\to N$, one obtains a representation of 
$f^!A$ on the (usual) pullback bundle $f^*E$. 

In this note, we only consider the case that $f$ is given by a surjective submersion 
%
\[\pi\colon M\to N.\]
%
In this case, transversality is automatic, and the Lie algebroid morphism $\pi^!A\to A$ is fiberwise surjective,  with kernel
$T_\pi M=\ker(T\pi)$ the tangent bundle to the foliation defined by $\pi$. Sections of $\wedge T_\pi^* M$ are called foliated (or leafwise) 
differential forms; we use the notation $\d_\pi$ for the Lie algebroid differential in order to avoid confusion with the de Rham differential.

\begin{theorem}[Crainic \cite{cra:dif}]\label{th:main}
Let $A\Ra N$ be a Lie algebroid, with a representation on a vector bundle $E\to N$. Suppose 	
$\pi\colon M\to N$ is a surjective submersion with connected fibers, such that the cohomology of the fibers vanishes in degrees $1\le k\le n$. 
Then the pullback map  
\[ \pi^*\colon H^\bullet(A,E)\to H^\bullet(\pi^!A,\pi^*E)\]
is an isomorphism in degrees $0\le k\le n$, and is injective in degree $n+1$. 	
\end{theorem}
%
As a special case of Theorem \ref{th:main}, taking $A=0_N$ to be the zero Lie algebroid with the trivial representation on $E=N\times \R$, we have that
$\pi^!A=T_\pi M$. The corresponding Lie algebroid cohomology is identified with the leafwise (foliated) cohomology, 
 \[ H^\bullet(T_\pi M)=H^\bullet_\pi(M).\]
One hence obtains: 
%
\begin{proposition}\label{prop:keycorollary}
Let $\pi\colon M\to N$ be a surjective submersion with connected fibers, such that the cohomology of the fibers vanishes in degrees $1\le k\le n$. 
Then the foliated cohomology $H^\bullet_\pi(M)$ vanishes in degrees  $1\le k\le n$.
\end{proposition}
%

\begin{remark}
The corresponding statement is false for general foliations $\F$: In particular, $H^1_\F(M)$ may be non-zero even if all leaves are simply connected. For example, let  $M=\R^2/\Z^2$ be the 2-torus, and let $\theta_1,\theta_2\in \Omega^1(M)$ be closed 1-forms pulling back to 
$\d x$ and $\d y+\mu \d x$, where $\mu\not\in\Q$. 
The Kronecker foliation $\mathcal F$ given by
$\ker\theta_2\subseteq TM$ has all leaves diffeomorphic to
$\mathbb R$, and hence $H^1(L)=0$ for every leaf $L$.  The restriction of $\theta_1$ to leaves is a foliated form $\alpha\in \Omega^1_F(M)$. Its class in  $H^1_\F(M)$ is non-zero: For if $\alpha=\d_\F f$ with $f\in C^\infty(M)$, then $\d f-\theta_1$ vanishes on $T_\F M$ and hence is a  function times $\theta_2$. But then  $\int \theta_1\wedge\theta_2=\int (\d f)\wedge\theta_2=\int \d(f\theta_2)=0$, 
which is a contradiction since $\theta_1\wedge\theta_2$ is a volume form. 
\end{remark}


Similar to \cite{cra:dif}, we will prove Proposition \ref{prop:keycorollary} as a first step towards  Theorem \ref{th:main}. To highlight the difficulty, let us point out that the proposition is greatly simplified when $\pi$ is a locally trivial fibration whose fibers are manifolds of finite type. Indeed, this then 
reduces further to the case of direct products, and one may use the following fact:

\begin{proposition}[Products] \label{prop:product}
	Suppose $M=Q\times N$ 	is a product of two manifolds, where $Q$ is of finite type, and let $\pi\colon M\to N$ be the projection to the second factor. 
	Then 
	\[ H_\pi^\bullet(M)=H^\bullet(Q)\otimes C^\infty(N).\]
\end{proposition}
\begin{proof}
	We want to show that the map 
	%
	\begin{equation}\label{eq:themap}
		\iota\colon\	\Omega^\bullet(Q)\otimes C^\infty(N)\to \Omega^\bullet_\pi(Q\times N)
	\end{equation}
	induces an isomorphism in cohomology.
	If $Q$ is smoothly contractible, this follows by  the Poincar\'{e} lemma on $Q$, treating the dependence on $N$ as parameters. For arbitrary $Q$ of finite type, pick a finite good open cover of $Q$, and apply the usual Mayer-Vietoris argument \cite{bo:di}. 
\end{proof}





\section{The Buchdahl argument}\label{app:buch}
%
Throughout the following discussion, $\pi\colon M\to N$ denotes a surjective submersion. 
%
\begin{proposition}[Buchdahl \cite{buc:rel}]\label{prop:buch}
	Suppose that for some $n\ge 1$, we have that 
	\[H^n(\pi^{-1}(x))=0,\ \ H^{n-1}(\pi^{-1}(x))=0
	\]
	for all $x\in N$. (For $n=1$, replace this with the assumption that $H^1(\pi^{-1}(x))=0$ for all $x\in N$ and  the fibers are connected.) 
	Then 
	\[H^n_\pi(M)=0.\]  
\end{proposition}
%
We shall give a proof in Section \ref{subsec:buch} below, using a modified version of the argument in \cite{buc:rel} and elementary results from de Rham theory.  Clearly, Proposition \ref{prop:keycorollary} follows from Proposition \ref{prop:buch}. 

\subsection{Preparatory lemmas}
Recall \cite{bo:di} that a \emph{good cover}  of a manifold $M$ is an open
cover such that every non-empty finite intersection of members of the
cover is contractible. It is a standard fact \cite{gr:co1} that every manifold $M$ admits 
a countable good cover. A manifold is of \emph{finite type} if it admits a finite good cover. This then implies that the de Rham cohomology of $M$ is finite dimensional. 
Given a manifold $M$, not necessarily of finite type, there exists an exhaustion by open subsets $U_1\subset U_2\subset \ldots$ with compact closures, such that 
\[ \bigcup_a U_a=M, \ \ \ \ \ \ol{U}_a\subset U_{a+1},\] 
and each $U_a$ is of finite type. (Just take a good open cover by open subsets with compact closures, and take $U_a$ to 
consist of suitable increasing finite unions of that cover.) Given such an exhaustion, we have: 

\begin{lemma}
	Suppose $H^k(M)=0$ for some $k$. Then there exists, for every $a$, some $b>a$ such that 
	the map  $H^k(U_{b})\to H^k(U_{a})$ is the zero map.
\end{lemma}
\begin{proof}
	%We will argue using Poincar\'{e} duality,  since this is natural with respect to open inclusions.
	Let $H^\bullet_{\on{comp}}(M,o_M)$ be the cohomology with compact support, with values in the orientation bundle
	$o_M$.
	Poincar\'{e} duality \cite[page 194]{gr:co1}
	states that for every manifold $M$ (not necessarily of finite type), the natural map 
	\[ H^\bullet(M)\to 	(H_{\on{comp}}^{\dim M-\bullet}(M,o_M))^*\]
	is an isomorphism. The statement is hence equivalent to saying that 
	%In particular, $H^k(M)=0$ is equivalent to $H^\ell_{\on{comp}}(M,o_M)=0$, where  $\ell=\dim M-k$. It suffices to prove that 
	if	$H^{\ell}_{\on{comp}}(M,o_M)=0$ for some $\ell$, then there exists, for every given $a$, some $b>a$ such that 
	the map  $H^\ell_{\on{comp}}(U_{a},o_M)\to H^\ell_{\on{comp}}(U_{b},o_M)$ is the zero map. Since 
	$U_a$ is of finite type, the cohomology group $H^\ell_{\on{comp}}(U_{a},o_M)$ is finite dimensional. We may therefore choose closed forms 
	\[ \alpha_1,\ldots,\alpha_r\in  \Omega^\ell_{\on{comp}}(M,o_M)
	\] 
	supported in $U_a$, such that their cohomology classes form a basis. Since $H^{\ell}_{\on{comp}}(M,o_M)=0$, 
	these forms admit primitives
	\[\beta_1,\ldots,\beta_r
	\in  \Omega^{\ell-1}_{\on{comp}}(M,o_M).\] 
	These primitives need not be supported in $U_a$, but we can choose  $b>a$ so that 
	$U_b$ contains the support of these primitives. It then follows that $\alpha_i|_{U_b}$ are exact.  
	Hence 
	$H^{\ell}_{\on{comp}}(U_a,o_M)\to  H^{\ell}_{\on{comp}}(U_b,o_M)$ is the zero map, and so is the dual map. 
\end{proof}

We shall need the analogous results for submersions  $\pi\colon M\to N$. 
%
\begin{lemma}[Fiberwise good covers]
	Given a submersion $\pi\colon M\to N$, we may choose a locally finite cover of $M$ by open sets, with compact closures, such that for every $x\in N$ the resulting cover of the fiber $\pi^{-1}(x)$ obtained by intersections is a good cover.
\end{lemma}
\begin{proof}
	Choose a Euclidean inner product on $T_\pi M=\ker(T\pi)$. This determines a family of Riemannian metrics on the fibers of $\pi$, depending smoothly on the base point in $N$.  
	For every $p\in M$, there exists an open neighborhood $U$ of $p$ with the 
	property that all intersections $U\cap \pi^{-1}(x)$ are geodesically convex, in the strong sense that any two points of this set are connected by a unique length minimizing geodesic segment inside this set. 
		(For example, we may choose a section of $\pi\colon M\to N$ over 
	an open subset $V\subset N$
	passing through $p$. After shrinking $V$ if needed, choose $\epsilon>0$ smaller than the fiberwise convexity radius along the image of the section. The fiberwise $\epsilon$-balls about the section then have convex intersections with every fiber.)
	 The intersection of geodesically convex subsets is again geodesically convex. A compact subset of $M$ is covered by finitely many subsets of this kind. 
	By the usual exhaustion argument, this leads to a locally finite cover of $M$ by  subsets of this kind. 
\end{proof}

\begin{lemma}[Fiberwise exhaustions]
	Let $\pi\colon M\to N$ be a submersion. Then there exists a nested sequence of open subsets $U_1\subset U_2\subset \cdots $ of $M$, with compact closures, 
	such that 
	\[ \bigcup_a U_a=M, \ \ \ \ \ \ol{U}_a\subset U_{a+1},\] 
	and such that each $\pi^{-1}(x)\cap U_a$ is of finite type for every $x\in N$. 
\end{lemma} 
\begin{proof}
	As in the case of $N=\pt$, we may build the $U_a$ by taking increasing finite unions of members of a fiberwise good cover. 
\end{proof}


\begin{lemma}[Product neighborhoods]
Given $x\in N$ and any open subset $U\subset M$ with compact closure, there exists an open neighborhood $V$ of $x$ and an embedding 
\[ 	\Phi\colon (U\cap \pi^{-1}(x))\times V\to M\]
intertwining projection to the second factor with the map $\pi$. 
\end{lemma}
\begin{proof}
	Such a neighborhood may be constructed using Euler-like vector fields \cite{bur:spl}. 
	The map $\pi$ induces a map of normal bundles 	$ \nu(M,\pi^{-1}(x))\to \nu(N,\{x\})=T_xN$, 
	which is a fiberwise isomorphism (since $\pi$ is a submersion). One hence has a trivialization
	\[
	\nu(M,\pi^{-1}(x))\cong \pi^{-1}(x)\times T_xN.
	\]
	Pick an Euler-like vector field on $N$ with respect to $\{x\}$. It admits a lift to an Euler-like vector field on $M$ with respect to $\pi^{-1}(x)$. One obtains open neighborhoods $O_M\subset \nu(M,\pi^{-1}(x))$ of $\pi^{-1}(x)$ 
	and $O_N\subset \nu(N,\{x\})$ of $\{x\}$,  with tubular neighborhood embeddings 
	\[ \phi_M\colon O_M\to M,\ \ \phi_N\colon O_N\to N\]
	which intertwine the submersion  $O_M\to O_N$ (given by restriction of $\nu(M,\pi^{-1}(x))\to  T_xN$) 
	with the given submersion $\pi$. The subset 
	$O_M$ may be strictly smaller than	$ \pi^{-1}(x)\times O_N$ in general. However, since $\ol{U}\cap \pi^{-1}(x)$ is compact, we 
	may choose an open neighborhood $V$ of $x$ such that $(\ol{U}\cap \pi^{-1}(x))\times \phi_N^{-1}(V)\subset O_M$. The map $\Phi$ is obtained by restriction of $\phi_M\circ (\on{id}\times \phi_N^{-1})$.
\end{proof}



\begin{lemma}\label{lem:3.6}
	Suppose that for some $k$, we have that 
	$H^k(\pi^{-1}(x))=0$ for all $x\in N$. Given an exhaustion as in the previous lemma, there exists, for all $a$, some $b>a$ such that the restriction map  
	\[ H^k_\pi(U_{b})\to H^k_\pi(U_a)\]
	is zero. 
\end{lemma}


\begin{proof}\phantom{.}
%\begin{enumerate}\renewcommand{\theenumi}{\arabic{enumi}}	\renewcommand{\labelenumi}{\theenumi.}
	%\item 	
	Given $x\in N$, fix  $c>a+1$ so that 	
	\[ H^k(U_{c}\cap \pi^{-1}(x))\to H^k(U_{a+1}\cap \pi^{-1}(x))\]
	is the zero map. 
	%\item 
	Choose a product open neighborhood $\Phi((U_c\cap \pi^{-1}(x))\times V)$ as in 
	the previous lemma. 
	%\item 
	Shrinking $V$ if necessary, the product neighborhood  $\Phi((U_{a+1}\cap \pi^{-1}(x))\times V)$ contains 
	$\ol{U}_{a}\cap \pi^{-1}(V)$, and taking $b$ sufficiently large the neighborhood $\Phi((U_c\cap \pi^{-1}(x))\times V)$ 
	is contained in $U_b$. This gives inclusions
	\[ U_{a}\cap \pi^{-1}(V)\hra  \Phi((U_{a+1}\cap \pi^{-1}(x))\times V)\hra \Phi((U_c\cap \pi^{-1}(x))\times V)
	\hra U_b\cap \pi^{-1}(V)
	.\]
   %\item 
  The map 
   \[ H^k_\pi (\Phi((U_c\cap \pi^{-1}(x))\times V))\to  H^k_\pi (\Phi((U_{a+1}\cap \pi^{-1}(x))\times V)),\]
 induced by the middle inclusion is the zero map, since 
 \[ H^k_\pi (\Phi((U\cap \pi^{-1}(x))\times V))=H^k(U\cap \pi^{-1}(x))\otimes C^\infty(V)\]
 by Proposition \ref{prop:product}. It follows that the map 
   %
   \[ H^k_\pi(U_{b}\cap \pi^{-1}(V))\to H^k_\pi(U_a\cap \pi^{-1}(V))\]
 is the zero map.  The base manifold $N$ may be covered by open subsets $V$ of this sort. Since 
   $U_a$ has compact closure, there is a finite collection $\{V_i\}$ of such open subsets, such that 
   $\{\pi^{-1}(V_i)\}$ cover $\ol{U}_a$. Choose $b_i$ such that the maps 
     \[ H^k_\pi(U_{b_i}\cap \pi^{-1}(V_i))\to H^k_\pi(U_a\cap \pi^{-1}(V_i))\]
  are all zero, and let $b$ be their maximum. Given $\alpha\in \Omega^k_\pi(U_b)$ with $\d_\pi\alpha=0$, 
  the restriction to $U_{a}\cap \pi^{-1}(V_i)$ admits a primitive $\beta_i$. 
  Let  $\{\chi_i\}$ be a partition of unity  subordinate to $\{V_i\}$ (regarded as a cover of $\bigcup V_i\subset N$) 
    and put $\beta=\sum_i (\pi^*\chi_i)\beta_i$. 
  Then $\d_\pi\beta=\alpha$ on $U_a$. 
\end{proof}

\subsection{Proof of Proposition \ref{prop:buch}} \label{subsec:buch}
Consider first the case $n=1$, thus assume that all fibers are connected and have vanishing first cohomology. 
Choose an open cover $\{V_i\}$ of $N$ such that $\pi$ admits a
section
\[s_i\colon V_i\to \pi^{-1}(V_i)\subset  M,\]
and let $\{\chi_i\}$ be a partition of unity subordinate to this cover. Given a $\d_\pi$-closed form 
 $\alpha\in\Omega^1_\pi(M)$  we obtain unique functions $f_i\colon \pi^{-1}(V_i)\to \R$ such that $f_i\circ s_i=0$ and 
the restriction of $f_i$ to each fiber $\pi^{-1}(x),\ x\in V_i$ is a primitive of the restriction of  $\alpha$. The function $f_i$ is smooth. (To see this near a given point $p\in \pi^{-1}(x)$, choose a path in $\pi^{-1}(x)$  from $p_0=s_i(x)$ to $p$. 
Subdivide the path into finitely many segments joining successive
points
\[
p_0,p_1,\ldots,p_N=p,
\]
 each of which is contained in a submersion chart. It then follows inductively that $f_i$ is smooth near $p_1$, hence 
smooth near $p_2$, and so on.)  The function 
\[
f=\sum_i (\pi^*\chi_i)f_i
\]
satisfies $\d_\pi f=\alpha$. This proves $H^1_\pi(M)=0$.

Consider now the case $n>1$. Lemma \ref{lem:3.6}, after passing to a subsequence, allows us to construct a sequence 
\[ U_1\subset U_2\subset \cdots \]
of open subsets of $M$, with compact closures, satisfying 
\[ \ol{U}_a\subset U_{a+1},\ \ \ \bigcup_a U_a=M,\] 
such that for all $a$ the maps 
	\[ H^n_\pi(U_{a+1})\to H^n_\pi(U_a),\ \ \ H^{n-1}_\pi(U_{a+1})\to H^{n-1}_\pi(U_a),
	\]
are zero. Let us also choose smooth functions $\lambda_a$ supported on $U_{a+1}$, with $\lambda_a|_{U_a}=1$. 

To show that $H^n_\pi(M)=0$, let a closed foliated form
\[ \alpha\in \Omega^n_\pi(M),\ \ \ \d_\pi\alpha=0\] 
be given. For every $a$, the restriction of $\alpha$ to $U_{a+3}$ is closed,
and hence its restriction to $U_{a+2}$ is exact. Choose
\[
\beta'_a\in\Omega^{n-1}_\pi(U_{a+2}),
\qquad
d_\pi\beta'_a=\alpha|_{U_{a+2}}.
\]
We shall modify these primitives inductively. Put
$\beta_1=\beta'_1$, and suppose that
\[
\beta_a\in\Omega^{n-1}_\pi(U_{a+2})
\]
has been constructed.  The difference $\beta'_{a+1}|_{U_{a+2}}-\beta_a\in \Omega^{n-1}_\pi(U_{a+2})$ is closed, and becomes exact after further restriction to $U_{a+1}$. Let 	
\[ \gamma_a\in \Omega_\pi^{n-2}(U_{a+1})\] 
be a primitive. Since $\lambda_a\gamma_a$ has compact support in $U_{a+1}$, it
extends by zero to $U_{a+3}$. Define
\[ \beta_{a+1}=\beta'_{a+1}-\d_\pi(\lambda_a \gamma_a).\]
Then $\beta_{a+1}$ agrees with $\beta_a$ on $U_a$, and is still a primitive of $\alpha|_{U_{a+3}}$. 
Consequently, there is a global form $\beta\in \Omega^{n-1}_\pi(M)$ with $\beta|_{U_a}=\beta_a$. 
This form  satisfies $\d_\pi\beta=\alpha$.  \qedhere

\subsection{Coefficients}
%
As already mentioned, Proposition \ref{prop:keycorollary} 
is a direct consequence of Proposition \ref{prop:buch}. In fact, we obtain a more general statement by considering coefficients. 
Let 
\[ E\to N\]
be a given vector bundle. Then $\pi^*E$ is fiberwise trivial, and has a canonical $T_\pi M$-representation, 
defining $H^\bullet_\pi(M,\pi^*E)$.

\begin{corollary}\label{cor:buchcorollary}
	Under the assumptions of Proposition \ref{prop:keycorollary}, let $E\to N$ be a vector
	bundle. Then
	\[
	H^k_\pi(M,\pi^*E)=0,
	\qquad 1\leq k\leq n,
	\]
	and
	\[
	H^0_\pi(M,\pi^*E)=\Gamma(E).
	\]
\end{corollary}
\begin{proof}
	The degree-zero statement follows from connectedness of the fibers:
	a smooth section of $\pi^*E$ that is constant along the fibers
	descends to a section of $E$. For the positive-degree
	statement, choose an open cover of $N$ over which $E$ is trivial.
	The assertion holds over the inverse image of each member of this cover by
	Proposition \ref{prop:buch}, applied componentwise. The local primitives may be
	patched using a partition of unity on $N$.
\end{proof}

\section{The pullback theorem and its extensions}
We now use the vanishing result from the previous section to prove
Theorem \ref{th:main}. After choosing a splitting of the pullback Lie algebroid,
the proof is a standard zig-zag argument. We then discuss some special cases and
extensions.

\subsection{Proof of Theorem \ref{th:main}}
We assume that the fibers of $\pi\colon M\to N$ are connected and have vanishing cohomology in degrees $1\le k\le n$. Let 
$A\Ra N$ be a Lie algebroid, and $E\to N$ a representation of $A$. The Lie algebroid complex of $A$ with coefficients in $E$ is denoted 
\[ \CC^\bullet(A,E)=\Gamma(\wedge^\bullet A^*\otimes E)\]
with the Chevalley-Eilenberg differential $\d_A$. 
Similarly, we have the Lie algebroid complex of $\pi^!A$ with coefficients in $\pi^*E$. The pullback map  
\begin{equation}\label{eq:pistarforms}
\pi^*\colon \CC^\bullet(A,E)\to \CC^\bullet(\pi^!A,\pi^*E)
\end{equation}
is a cochain map. After choosing a splitting of the exact sequence of vector bundles
\[ 0\to T_\pi M\to \pi^!A\to \pi^*A\to 0,\]
we have $\CC^k(\pi^!A,\pi^*E)=\bigoplus_{p+q=k}\CC^{p,q}$ with 
%
\[ \CC^{p,q}=\Omega^q_\pi(M,\pi^*(\wedge^p A^*\otimes E)).\]
Using the Chevalley-Eilenberg formula for the Lie algebroid differential $\d_{\pi^!A}$ and the decomposition 
$\pi^!A=T_\pi M\oplus \pi^*A$, we see that $\d_{\pi^!A}$ is the sum of 
%
\[ \d_\pi\colon \CC^{p,q}\to \CC^{p,q+1},\]
%
and two other summands of bidegrees $(1,0)$ and $(2,-1)$.  
By Corollary \ref{cor:buchcorollary}, the cohomology of $\CC^{p,\bullet}$ vanishes in degree $1\le q\le n$, and is given by 
$\Gamma(\wedge^p A^*\otimes E)$ for $q=0$. 

We now use the standard zig-zag argument: To show that the map 
\begin{equation}\label{eq:pistarcoh}
	 \pi^*\colon H^k(A,E)\to H^k(\pi^!A,\pi^*E)
\end{equation}
is injective for $k\le n+1$, let $\alpha\in \CC^k(A,E)$ be a $\d_A$-cocycle such that $\pi^*\alpha$ is exact. We need to show that $\alpha$ is exact. Write $\pi^*\alpha=\d_{\pi^! A}\beta$ for $\beta\in \CC^{k-1}(\pi^!A,\pi^*E)$. 
It decomposes as 
\[ \beta=\beta_q+\ldots+\beta_0\]
with $\beta_i\in \CC^{k-1-i,i}$. Then $\d_\pi \beta_q=0$. If $q>0$, it follows that $\beta_q=\d_\pi \gamma$
for some element $\gamma\in \CC^{k-1-q,q-1}$. 
Replacing $\beta$ with $\beta-\d_{\pi^!A}\gamma$, we have arranged that $\beta_q=0$, but $\d_\pi \beta_{q-1}=0$. 
Proceeding in this manner, 
we eventually arrive at an element  $\beta\in \CC^{k-1,0}$ with $\d_{\pi^!A}\beta=\pi^*\alpha$. This means in particular $\d_\pi\beta=0$, hence 
$\beta=\pi^*\gamma$ for some $\gamma\in \CC^{k-1}(A,E)$. It follows that 
\[ \pi^*\d_A\gamma=\d_{\pi^!A}\beta=\pi^*\alpha,\]
hence $\d_A\gamma=\alpha$. This shows $\alpha$ is exact, as required. 
Similarly, to see that the map \eqref{eq:pistarcoh} is surjective for $k\le n$, let $\sigma\in  \CC^{k}(\pi^!A,\pi^*E)$ be a cocycle
representing a given cohomology class. Write 
$\sigma=\sigma_q+\ldots+\sigma_0$ with $\sigma_i\in \CC^{k-i,i}$. Then $\d_\pi\sigma_q=0$. If $q>0$, it follows that 
$\sigma_q=\d_\pi\tau$ for some $\tau\in \CC^{k-q,q-1}$. Replacing $\sigma$ with $\sigma-\d_{\pi^!A}\tau$, we have arranged $\sigma^q=0$. 
Proceeding in this manner, we see that the given cohomology class may be represented by an element $\sigma$ in $\CC^{k,0}$. Since 
$\d_\pi\sigma=0$, this element is of the form $\sigma=\pi^*\alpha$ for some $\alpha\in \CC^k(A,E)$, with $\pi^*\d_A\alpha
=\d_{\pi^!A}\pi^*\alpha=\d_{\pi^!A}\sigma=0$, hence $\d_A\alpha=0$. 
 \qed
 
 
 \subsection{Special cases and remarks}
 
 \subsubsection{The case $A=TN$}
 We have already mentioned that for the zero Lie algebroid $A=0_N$, with $E=N\times \R$, Theorem \ref{th:main} reduces to Proposition \ref{prop:keycorollary}. Another noteworthy setting is $A=TN$, so that $\pi^!A=TM$. An $A$-representation on a vector bundle $E\to N$ is the same as a flat connection on $E$, and $H^\bullet(A,E)$ is the de Rham cohomology with coefficients in the flat vector bundle $E$. 
 Theorem \ref{th:main}  says that if $\pi\colon M\to N$ has connected fibers with vanishing cohomology in degrees $1\le k\le n$, then the pullback map 
 \[ H^k(N,E)\to H^k(M,\pi^*E)\] 
  is an isomorphism in degrees $0\le k\le n$, and is injective for $k=n+1$. 
  
 \subsubsection{Complex Lie algebroids} 
 	Throughout our discussion, it was tacitly assumed that $A$ is a \emph{real} Lie algebroid. The argument also applies to \emph{complex} Lie algebroids (with anchor map  taking values in the complexified tangent bundle $T^\C N$), with representations on complex vector bundles $E\to N$. 
 	Here the pullback $\pi^!A$ is defined as the fiber product 
 	$\T^\C M\times_{T^\C N} A$, and comes with a Lie algebroid morphism to $A$ with kernel $T_\pi^\C M$. The statement and proof of Theorem \ref{th:main} extend to this setting, keeping in mind that $\CC^\bullet (A,E)$ and $\CC^\bullet(\pi^!A,\pi^*E)$ are now complex vector spaces. 
 
 
 \subsubsection{Non-Hausdorff bases} 
 Recall that by our standing assumption, all manifolds in this note are required to be Hausdorff. 
 Indeed, the conclusions of Theorem \ref{th:main} and Proposition \ref{prop:keycorollary}
 may fail without this assumption. As an example, let 
 \[
 M=\mathbb R^2\setminus\{(0,0)\}
 \]
 with the foliation given by connected components of the vertical lines. The leaf space $N$ is the real line with doubled origin, and  is a standard example of a non-Hausdorff manifold. The quotient map $\pi\colon M\to N$ is a surjective submersion, with contractible fibers. 
 Nevertheless, $H^1_\pi(M)\neq 0$. A non-trivial class is represented by $\alpha\in \Omega^1_\pi(M)$, given by 
restriction of the angular form 
 \[
 \omega=\frac{x \d y-y\d x}{x^2+y^2}\in\Omega^1(M)
 \]
 to the leaves. Suppose $\alpha=d_\pi g$ for a function $g\in C^\infty(M)$. Then $\omega-\d g$ vanishes on $T_\pi M$, and hence is of the form 
 %
 \[ \omega-\d g=\phi(x,y)\d x\] 
 for some function $\phi\in C^\infty(M)$. Let $\psi(x)=\phi(x,1)$. 
 Since $\omega$ is closed, $\phi$ must be independent of $y$, and so $\phi(x,y)=\psi(x)$ for all $y$. 
 (By continuity, this even holds true for the disconnected line $x=0$.) 
 We find that $\omega=\d g+\psi(x)\d x$ is exact,  contradicting the fact that  $\omega$ represents a generator of $H^1(M)=\R$. 
 
 
 
 
\subsection{The holomorphic setting}

Let $A \Rightarrow N$ be a holomorphic Lie algebroid, with a
representation on a holomorphic vector bundle $E \to N$. Following \cite[Section 6]{ev:tra}, 
one defines the
holomorphic Lie algebroid cohomology 
\[ H^\bullet_{\mathrm{hol}}(A,E)\] 
as the hypercohomology of the complex of sheaves $\mathcal O(\wedge^\bullet A^* \otimes E)$ 
with the Chevalley-Eilenberg differential.  It may alternatively be computed as the cohomology of the total complex of the 
Dolbeault double complex
\[
\Omega^{0,s}
\bigl(
N,\wedge^r A^* \otimes E
\bigr),
\]
with differentials $d_A$ and $\bar\partial$, of bidegrees $(1,0)$ and
$(0,1)$. Laurent-Gengoux, Sti\'enon, and Xu \cite{lau:hol} showed that $T^{0,1}N$ and
$A^{1,0}$ form a matched pair of complex Lie algebroids, and hence combine into a complex Lie algebroid
\[
A_{\mathrm{Dol}}
=
T^{0,1}N \bowtie A^{1,0}
\]
which we shall call the \emph{Dolbeault Lie algebroid}. 
The holomorphic vector bundle $E$ carries a natural representation of $A_{\mathrm{Dol}}$, and its
Chevalley--Eilenberg complex is naturally identified with the total
Dolbeault complex above. Consequently, by \cite[Theorem 4.19]{lau:hol},
%
\begin{equation}\label{eq:gsx}
	H^\bullet_{\mathrm{hol}}(A,E)
	\cong
	H^\bullet(A_{\mathrm{Dol}},E).
\end{equation}
%
We next consider the functoriality of this construction. 

\begin{lemma}\label{lem:functorial}
	Let $f\colon M\to N$ be a holomorphic map, transverse to the anchor
	of the holomorphic Lie algebroid $A\Rightarrow N$. Then there is a
	canonical isomorphism of complex Lie algebroids
	\[
	f^!(A_{\mathrm{Dol}})\cong (f^!A)_{\mathrm{Dol}},
	\]
	compatible with the induced representations on
	$f^*E$. The following diagram commutes:
	\[
	\xymatrix@C=4.5em@R=2.5em{
		H^\bullet_{\mathrm{hol}}(A,E)
		\ar[r]^{f^*}
		\ar[d]_{\cong}
		&
		H^\bullet_{\mathrm{hol}}(f^!A,f^*E)
		\ar[d]^{\cong}
		\\
		H^\bullet(A_{\mathrm{Dol}},E)
		\ar[r]_{f^*}
		&
		H^\bullet((f^!A)_{\mathrm{Dol}},f^*E)
	}
	\]
\end{lemma}
\begin{proof}
	Let 
\[ 		(v,(w,a))\in  f^!(A_{\mathrm{Dol}})\subset T_\C M\times f^*(T^{0,1}N\oplus A^{1,0}).\]
Thus $T_{\mathbb C}f(v)=w+\varrho(a)$. 
Decomposing into holomorphic and anti-holomorphic parts, and using that 
$f$ is holomorphic, we have that
	\[
	w=(T^{0,1}f)(v^{0,1}),
	\qquad
	(T^{1,0}f)(v^{1,0})=\varrho(a).
	\]
	Consequently, the map taking $(v,(w,a))$ to $\bigl(v^{0,1},(v^{1,0},a)\bigr)$ defines an isomorphism of complex  
	vector bundles
	\[
	f^!(A_{\mathrm{Dol}})
	\to 
	T^{0,1}M\oplus(f^!A)^{1,0}
	=(f^!A)_{\mathrm{Dol}}.
	\]
	The definitions of the Dolbeault Lie algebroid structures show that this is an isomorphism of
	complex Lie algebroids. It also intertwines the induced
	representations on $f^*E$.
	
	Under the identification of the Chevalley--Eilenberg complexes of the
	Dolbeault Lie algebroids with the corresponding total Dolbeault
	complexes, the induced cochain maps are both given by pullback.
	The diagram therefore commutes.
\end{proof}




Specializing to the case of submersions, we arrive at the following result. 

\begin{theorem}
Let $\pi\colon M\to N$ be a holomorphic submersion with connected
fibers whose ordinary de Rham cohomology vanishes in degrees
$1\leq k\leq n$.
Let $A\Ra N$ be a holomorphic Lie algebroid, with a representation on a holomorphic vector bundle $E\to N$. Then 
the pullback map 
\[ \pi^*\colon H^\bullet_{\mathrm{hol}}(A,E)\to H^\bullet_{\mathrm{hol}}(\pi^!A,\pi^*E)\]
is an isomorphism in degrees $0\le k\le n$, and is injective in degree $n+1$. 
\end{theorem}

\begin{proof}
Using the isomorphism \eqref{eq:gsx}, together with the functoriality statement from Lemma \ref{lem:functorial}, this 
is equivalent to the statement that $\pi^*\colon H^\bullet(A_{\on{Dol}},E)\to  H^\bullet(\pi^!(A_{\on{Dol}}),\pi^*E)$ 
is an isomorphism in these degrees. But this is just Theorem \ref{th:main} for the case of complex Lie algebroids. 
\end{proof}


%\bibliographystyle{amsplain-initials}
%\bibliography{../../../../Biblio/ref}

\bibliographystyle{amsplain} 
\def\cprime{$'$} \def\polhk#1{\setbox0=\hbox{#1}{\ooalign{\hidewidth
			\lower1.5ex\hbox{`}\hidewidth\crcr\unhbox0}}} \def\cprime{$'$}
\def\cprime{$'$} \def\cprime{$'$} \def\cprime{$'$} \def\cprime{$'$}
\def\polhk#1{\setbox0=\hbox{#1}{\ooalign{\hidewidth
			\lower1.5ex\hbox{`}\hidewidth\crcr\unhbox0}}} \def\cprime{$'$}
\def\cprime{$'$} \def\cprime{$'$} \def\cprime{$'$} \def\cprime{$'$}
\providecommand{\bysame}{\leavevmode\hbox to3em{\hrulefill}\thinspace}
\providecommand{\MR}{\relax\ifhmode\unskip\space\fi MR }
% \MRhref is called by the amsart/book/proc definition of \MR.
\providecommand{\MRhref}[2]{%
	\href{http://www.ams.org/mathscinet-getitem?mr=#1}{#2}
}
\providecommand{\href}[2]{#2}
\begin{thebibliography}{1}
	
	\bibitem{ber:equ}
	Joseph Bernstein and Valery Lunts, \emph{Equivariant sheaves and functors},
	Lecture Notes in Mathematics, vol. 1578, Springer-Verlag, Berlin, 1994.
	
	\bibitem{bo:di}
	Raoul Bott and Loring Tu, \emph{Differential forms in algebraic topology},
	Graduate Texts in Mathematics, vol.~82, Springer-Verlag, New York, 1982.
	
	\bibitem{buc:rel}
	Nicholas~P. Buchdahl, \emph{{On the relative de Rham sequence}}, Proceedings of
	the American Mathematical Society \textbf{87} (1983), no.~2, 363--366.
	
	\bibitem{bur:spl}
	Henrique Bursztyn, Hudson Lima, and Eckhard Meinrenken, \emph{{Splitting
			theorems for Poisson and related structures}}, J. Reine Angew.~Math.~
	\textbf{754} (2019), 281--312.
	
	\bibitem{cra:dif}
	Marius Crainic, \emph{Differentiable and algebroid cohomology, van {E}st
		isomorphisms, and characteristic classes}, Comment. Math. Helv. \textbf{78}
	(2003), no.~4, 681--721.
	
	\bibitem{cra:lec}
	Marius Crainic, Rui~Loja Fernandes, and Ioan M{\u{a}}rcu{\c{t}}, \emph{Lectures
		on {P}oisson geometry}, Graduate Studies in Mathematics, vol. 217, American
	Mathematical Society, Providence, Rhode Island, 2021.
	
	\bibitem{ev:tra}
	Sam Evens, Jiang-Hua Lu, and Alan Weinstein, \emph{{Transverse measures, the
			modular class, and a cohomology pairing for Lie algebroids}}, The Quarterly
	Journal of Mathematics \textbf{50} (1999), no.~200, 417--436.
	
	\bibitem{gr:co1}
	Werner Greub, Stephen Halperin, and Ray Vanstone, \emph{{Connections,
			Curvature, and Cohomology. Volume I: De Rham Cohomology of Manifolds and
			Vector Bundles}}, Pure and Applied Mathematics, vol. 47-I, Academic Press,
	New York and London, 1972.
	
	\bibitem{lau:hol}
	Camille Laurent-Gengoux, Mathieu Sti\'{e}non, and Ping Xu, \emph{{Holomorphic
			Poisson Manifolds and Holomorphic Lie Algebroids}}, International Mathematics
	Research Notices \textbf{2008} (2008), rnn088, 46 pp.
	
\end{thebibliography}


\end{document}



\section{Further thoughts}
2. There should be a better formulation using colimits. Recall that 
\[ \on{colim}_{a\to \infty} H^n(U_a)\]
consists of sequences of classes in $H^n(U_a)$ that are related by the restriction morphisms. We have a global restriction map 
\[ H^n(M)\to \on{colim}_{a\to \infty} H^n(U_a),\]
%
and one may ask if this map is an isomorphism. The statement that, for all $a$ there exists $b>a$ such that the restriction is the zero map, means exactly that $\on{colim}_{a\to \infty} H^n(U_a)=0$. So, the first lemma says that if
$H^n(M)=0$ then  $\on{colim}_{a\to \infty} H^n(U_a)=0$. 

3. Can one prove a similar statement for regular Lie algebroids? That is: Given $A\Ra M$ a regular Lie algebroid, and assuming that 
for all leaves $L\subset M$ the Lie algebroid cohomology of $A_L$ vanishes in degrees $k,k-1$, then 
$H(A)$ vanishes in degree $k$?




\section{Foliations}
We now prove the analogous results for foliations $T_\F M\subset TM$. 

\begin{lemma}[Leafwise good covers]
	Given a foliation $\F$, there exists a locally finite cover of $M$ by open sets, with compact closures, such that for 
	every leaf $L\subset M$ the resulting cover of $L$ obtained by intersections is a good cover.
\end{lemma}
\begin{proof}
	Choose a Euclidean inner product on $T_\F M$. 
	This determines Riemannian metrics on the leaves.  For every $p\in M$, 
	there exists an open neighborhood $U$ with the 
	property that all intersections $U\cap L$ are geodesically convex. (For example, we may choose a local transversal $V$ through $p$, and then take points within $\epsilon>0$  of that transversal. For $\epsilon>0$ sufficiently small, and shrinking $V$ is needed, this will have the desired property.) A compact subset of $M$ is covered by finitely many open subsets of this kind. 
	This has the desired property, since the intersection of geodesically convex subsets is again geodesically convex. 
\end{proof}

\begin{remark}
	It appears that this lemma works equally well for foliations $\F$: There is a covering of $M$ by open subsets with the property that the intersection of the cover with any leaf is a good cover of that leaf. Such a cover may be constructed from any choice of a Riemannian metric on $T_\F M$.
\end{remark}


\begin{lemma}[Leafwise exhaustions]
	There exists a nested sequence of open subsets $U_a\subset M$, with compact closures, 
	such that 
	\[ \bigcup_a U_a=M, \ \ \ \ \ \ol{U}_a\subset U_{a+1},\] 
	and such that each $L\cap U_a$ is of finite type. 
\end{lemma} 
\begin{proof}
	As in the case of a trivial foliation, we may build the $U_a$ by taking increasing finite unions of members of a fiberwise good cover. 
\end{proof}

\begin{lemma}
	Let $A\Ra M$ be a regular Lie algebroid, and $\F$ the foliation defined by its anchor. 	Suppose that for some $k\in \N$ we have that $H^k(A_L)=0$ for all leaves. Given an exhaustion as in the previous lemma, there exists, for all $a$, some $b>a$ such that the restriction maps  
	\[ H^k(A|_{U_{b}})\to H^k(A|_{U_a})\]
	are all zero. 
\end{lemma}

\begin{proof}
	Begin by choosing $c>a$ so that 
	\[ H^k(A_L|_{L\cap U_{c}})\to H^k(A_L|_{L\cap U_a})\]
	is the zero map. Since $L\cap U_c$ has compact closure, there exists a tubular neighborhood [...] [...]
	
	
	Using a partition of unity on $N$, it suffices to show that every $x\in N$ has an open neighborhood $V$ such that 
	\[ H^n_\pi(U_{b}\cap \pi^{-1}(V))\to H^n_\pi(U_a\cap \pi^{-1}(V))\]
	is the zero map for $b$ sufficiently large. Begin by choosing $c>a$ so that 
	\[ H^n_\pi(U_{c}\cap \pi^{-1}(x))\to H^n_\pi(U_a\cap \pi^{-1}(x))\]
	is the zero map. Since $ U_{c}\cap \pi^{-1}(x)$ has compact closure, we may choose $V$ sufficiently small so that there is a tubular neighborhood embedding $(U_{c}\cap \pi^{-1}(x))\times V\to M$, intertwining $\pr_2$ with $\pi$. 
	Now choose $b$ sufficiently large so that $U_b$ contains  this tubular neighborhood. \footnote{Perhaps this strategy doesn't work so well, since the $A_L$'s for nearby leaves are not related in a simple way.} 
	
\end{proof}




\begin{remark}
	Question: Does the analogous fact hold for transitive Lie algebroids $A\Ra M$?  Evens-Lu-Weinstein establish a Poincar\'{e} duality pairing for Lie algebroids, which gives a map
	% 
	\[ H^\bullet(A)\to (H^\bullet_{\on{comp}}(A,Q_A\otimes o_M))^*       \]
	using the $A$-representation on the line bundle $Q_A=\wedge^{\on{top}}TA\otimes \wedge^{\on{top}}T^*M$. They point out that this map is not an isomorphism, in general. An unpublished paper of Kobarski-Michchenko 
	\url{https://www.impan.pl/~jagrab/kubarski.pdf} states that the pairing is non-degenerate when $A$ is transitive. Their proof is somewhat incomprehensible, but one would expect that the GHV argument can be adapted. But if we take the result for granted, 
	it seems that the exact same argument will show that if $H^k(A)=0$, then for all $a$ there exists $b>a$ such that $H^k(A|_{U_b})\to H^k(A|_{U_a})$ is the zero map. 
\end{remark}





\end{document}
